% SOURCE_STATEMENT: "Prove mean-field convergence for Keller–Segel-type and Dirac-approximating potentials, and obtain the better-than-$N^{-1/2}$ mean-square rate needed to derive intermediate fluctuation limits from smoothed empirical measures."
% Preparation and final self-audit: CONTRIBUTING.md prompts C and D; source/audit artifacts retained in the local physics_statements_2026_10_11 batch.
% !TeX program = lualatex
\documentclass[11pt]{article}
\usepackage[margin=25mm]{geometry}
\usepackage{fontspec}
\setmainfont{Latin Modern Roman}
\usepackage{amsmath,amssymb,mathtools}
\usepackage{unicode-math}
\setmathfont{Latin Modern Math}
\usepackage{xurl,hyperref}
\hypersetup{colorlinks=true,urlcolor=blue}
\setcounter{secnumdepth}{0}
\setlength{\parindent}{0pt}
\setlength{\parskip}{5pt}
\setlength{\emergencystretch}{3em}
\newcommand{\sref}[2]{\hyperlink{source:#1:#2}{[S#2]}}
\newcommand{\cataloguescope}[1]{\paragraph{Recorded target.}#1}
\newcommand{\R}{\mathbb R}
\newcommand{\C}{\mathbb C}
\newcommand{\N}{\mathbb N}
\newcommand{\Z}{\mathbb Z}
\newcommand{\Q}{\mathbb Q}
\newcommand{\F}{\mathbb F}
\newcommand{\eps}{\varepsilon}
\begin{document}
% UnsolvedMath ID 30005625; source catalogue code OWR-14297737-002
\setcounter{section}{30005624}
\section{Fluctuation Limits for Moderately Interacting Particle Systems}\label{30005625}
{\small UnsolvedMath ID 30005625 · Mathematical Physics}
\subsection{Definitions and mathematical statement}

\paragraph{Shared notation.}$\mathbb N=\{1,2,\ldots\}$; the standard number fields have their usual meanings. All additional conventions are stated locally.


For particle positions $X_i(t)$, the empirical measure is $\mu_N(t)=N^{-1}\sum_{i=1}^N\delta_{X_i(t)}$. The source uses $V_\eta=Z_\eta*Z_\eta$, a smooth normalized symmetric compactly supported convolution-square potential, with algebraic interaction length $\eta=N^{-\beta}$. Its intermediate nonlocal density $u^\eta$ solves the report's aggregation--diffusion equation. Smoothed measures are $f^{N,\eta}=\mu^{N,\eta}*Z_\eta$ and $g^\eta=u^\eta*Z_\eta$. Intermediate fluctuations are $\sqrt N(\mu^{N,\eta}-u^\eta)$, distinct from fluctuations around the final local limit $u$.

\paragraph{Catalogue wording retained for traceability.}
Prove mean-field convergence for Keller–Segel-type and Dirac-approximating potentials, and obtain the better-than-$N^{-1/2}$ mean-square rate needed to derive intermediate fluctuation limits from smoothed empirical measures.

\paragraph{Scope and source qualification.}The report already proves a better-than-$N^{-1/2}$ smoothed error estimate conditional on an assumed mean-field convergence result. The remaining task is to verify that premise for the specified additional potentials and derive the fluctuation limit. No interchangeable interpretation of mean-square error and its square root is assumed.
\subsection{Short English statement}
Obtain quantitative mean-field limits strong enough to control intermediate fluctuations for the stated singular or shrinking-range interaction potentials.
\subsection{Sources}
\begin{itemize}
\item[\hypertarget{source:30005625:1}{S1}] ULAM.AI and the UnsolvedMath Contributors, UnsolvedMath, record 30005625 (OWR-14297737-002). Catalogue curation under CC BY 4.0; definitions and English exposition added locally.
\url{https://huggingface.co/datasets/ulamai/UnsolvedMath}.
\item[\hypertarget{source:30005625:2}{S2}] Classical and Quantum Mechanical Models of Many-Particle Systems (2023), Oberwolfach Reports, DOI 10.4171/owr/2023/38, pp. 2179--2246. Original report and mathematical context.
\url{https://doi.org/10.4171/owr/2023/38}.
\end{itemize}
\label{end:30005625}
\end{document}
